% ============ 附录 A 关于群、环、域的一些知识 (原书 pp.403-414) ============
\chapter*{附录 A: 关于群、环、域的一些知识}

在这个附录中, 我们简要地叙述关于群、环、域的一些基本概念和事实. 除了一个例外, 我们略去了全部证明, 这个例外是: 我们在附录的最后给出代数基本定理（复数域的代数封闭性）的一个代数化的证明.

\section*{1. 有限生成 Abel 群}

设 $G$ 是 Abel（加法）群. 如果存在有限个元素 $g_1,\cdots,g_r\in G$, 使得 $G$ 中每个元素 $g$ 均可表示成
\[
  g=n_1g_1+\cdots+n_rg_r\qquad(n_i\in\Z,\ \Z\ \text{表示整数集合}),
\]
则群 $G$ 就叫作是有限生成的, 而 $g_1,\cdots,g_r$ 叫作是群 $G$ 的一组生成元素. $G$ 中有限阶元素 $g$（即存在 $n>0$, 使得 $ng=0$）叫作是扭元素. $G$ 中全部扭元素形成一个子群, 称作是 $G$ 的\emph{挠子群}, 表示成 $G_t$. 如果 $G_t=(0)$, 则 $G$ 叫作是无扭的. 如果 $G$ 是无扭群, 并且存在一组有限生成元素 $g_1,\cdots,g_r$, 使得 $G$ 中每个元素均可惟一地表示成公式 (1) 的形式, 则称 $G$ 是有限生成的自白 Abel 群. 换句话说, 有限生成自由 Abel 群 $G$ 是它的 $r$ 个子群 $\Z g_i=\{\Z g_i\}$ 的直和:
\[
  G=\Z g_1\oplus\cdots\oplus\Z g_r.
\]
这时, $g_1,\cdots,g_r$ 称作是自由 Abel 群 $G$ 的一组基, 而 $r$ 叫作是自由 Abel 群 $G$ 的秩, 表示成 $\operatorname{rank}G$.

\textbf{(D)（有限生成 Abel 群结构定理）}
(a) 若 $G$ 为 $n$ 阶有限 Abel 群, $n=p_1^{\alpha_1}\cdots p_s^{\alpha_s}$, 其中 $p_1,\cdots,p_s$ 是 $s$ 个不同的素数, $\alpha_i\geqslant 1$, 则 $G$ 为它的 $s$ 个 Sylow 子群 $G_{p_i}$ 的直和:
\[
  G=G_{p_1}\oplus\cdots\oplus G_{p_s}.
\]
(b) 每个有限生成 Abel 群 $G$ 均可表示成直和 $G=G_t\oplus G_f$, 其中 $G_t$ 是 $G$ 的扭子群（为有限群）, 而 $G_f$ 是 $G$ 的有限生成自由 Abel 子群.

(c) 有限生成自由 Abel 群 $G$ 的每个子群 $H$ 也是有限生成自由 Abel 群, 并且 $r=\operatorname{rank}H\leqslant\operatorname{rank}G=n$. 进而, 我们可以找到 $G$ 的一组基 $g_1,\cdots,g_n$, 使得
\[
  G=\Z g_1\oplus\cdots\oplus\Z g_n,\qquad
  H=\Z d_1g_1\oplus\cdots\oplus\Z d_rg_r,
\]
其中 $d_i\mid d_{i+1}\ (1\leqslant i<r)$, 并且 $d_i>0$.

\section*{2. 环的理想特性和元素特性}

本书中的环均指带 $1$ 交换环. 设 $S$ 是环 $R$ 的一个子集合. 定义
\[
  SR=\left\{\sum_{\text{有限}}r_is_i\mid r_i\in R,\ s_i\in S\right\},
\]
这是环 $R$ 的理想, 而且是环 $R$ 中包含集合 $S$ 的最小理想, 称作是由集合 $S$ 生成的理想. 如果 $S=\{a_1,\cdots,a_m\}$ 是 $R$ 的有限子集合, 则 $SR$ 也记作 $(a_1,\cdots,a_m)$, 并且称作是有限生成的理想. 特别当 $S=\{a\}$ 时, $(a)$ 叫作是环 $R$ 的主理想. 若环 $R$ 的每个理想都是主理想, 则称 $R$ 为主理想环. 设 $I_1,\cdots,I_m$ 均是环 $R$ 的理想, 定义
\[
  I_1+\cdots+I_m=\{a_1+\cdots+a_m\mid a_i\in I_i\},
\]
这也是环 $R$ 的理想, 叫作是理想 $I_1,\cdots,I_m$ 之和. 例如由集合 $\{a_1,\cdots,a_m\}$ 生成的理想即是 $Ra_1+\cdots+Ra_m$. 如果 $I_1+I_2=(1)=R$, 则称理想 $I_1$ 和 $I_2$ 互素.

环 $R$ 中元素 $a$ 叫作是 $R$ 中的零因子, 如果 $a\neq 0$ 并且存在 $0\neq b\in R$, 使得 $ab=0$. 没有零因子的（带 $1$ 交换）环叫作是整环. 如果整环 $D$ 同时是主理想环, 则称 $D$ 为主理想整环. 每个整环均可嵌在一个域中, 并且具有这样性质的最小域不计同构是惟一决定的, 称作是整环 $D$ 的商域. 环 $R$ 中的乘法可逆元叫作是 $R$ 中的单位. 环 $R$ 中全部单位形成乘法群, 叫作是环 $R$ 的单位群, 表示成 $U(R)$. 当 $R\neq 0$ 并且 $U(R)=R-\{0\}$ 时, $R$ 就叫作是域.

设 $\mathfrak{p}$ 是环 $R$ 的一个理想, $\mathfrak{p}\neq R$. 称 $\mathfrak{p}$ 是环 $R$ 的素理想, 是指: $\mathfrak{p}\neq R$, $a,b\in R$, $ab\in\mathfrak{p}\Rightarrow a\in\mathfrak{p}$ 或者 $b\in\mathfrak{p}$. 称 $\mathfrak{p}$ 是环 $R$ 的极大理想, 是指: $R$ 中不存在理想 $I$, 使得 $\mathfrak{p}\subsetneq I\subsetneq R$. 我们有

\textbf{(2)} (a) $\mathfrak{p}$ 为环 $R$ 的素理想 $\Leftrightarrow R/\mathfrak{p}$ 为整环;

(b) $\mathfrak{p}$ 为环 $R$ 的极大理想 $\Leftrightarrow R/\mathfrak{p}$ 为域.

由于域均是整环, 从而极大理想均是素理想. 利用 Zorn 引理可以证明: 任意（带 $1$ 交换）环中至少存在一个极大理想, 从而也至少存在一个素理想.

\textbf{(3)（中国剩余定理）} 设 $I_1,\cdots,I_n$ 是环 $R$ 中两两互素的理想, 则有环的自然同构
\[
  R/(I_1\cap I_2\cap\cdots\cap I_n)\cong R/I_1\times R/I_2\times\cdots\times R/I_n,
\]
\[
  a\ (\mathrm{mod}\ I_1\cap\cdots\cap I_n)\mapsto
  (a\ (\mathrm{mod}\ I_1),\cdots,a\ (\mathrm{mod}\ I_n)).
\]

设 $R$ 为环, $a,b\in R$, $b\neq 0$. 如果存在 $c\in R$, 使得 $a=bc$, 我们称 $a$ 整除 $b$, 表示成 $a\mid b$. 这时 $a$ 叫作是 $b$ 的因子, 而 $b$ 叫作是 $a$ 的倍元. 如果不存在满足上述条件的 $c$, 则称 $a$ 不能整除 $b$, 表示成 $a\nmid b$. 设 $a,b\in R'=R-\{0\}$. 如果 $a\mid b$ 并且 $b\mid a$, 则元素 $a$ 和 $b$ 叫作是相伴的, 表示成 $a\sim b$. 当 $R$ 是整环时, $\sim$ 是 $R$ 上的等价关系, 这是因为: $a\sim b\Leftrightarrow a$ 和 $b$ 相差一个单位因子. 假设 $a,b,c\in R$, $a=bc$, 并且 $b$ 和 $c$ 均不是 $R$ 中单位（即 $b$ 和 $c$ 均不与 $a$ 相伴）, 则 $b$ 和 $c$ 叫作是 $a$ 的真因子. $R'$ 中元素 $a$ 叫作是环 $R$ 的不可约元素, 如果 $a\notin U(R)$ 并且没有真因子. 环 $R$ 中这些元素特性与其理想特性的联系是:
\[
  a\in U(R)\Leftrightarrow(a)=R\ (\text{即 } (a)=Rr\ \text{对于每个 } r\in R).
\]
整环 $R$ 叫作是惟一因子分解整环（简记作 UFD）, 是指

(i)（分解的存在性）每个非零非单位元素 $a\in R$ 均可写成 $R$ 中有限个不可约元素之积;

(ii)（分解的惟一性）如果 $a=p_1\cdots p_m=q_1\cdots q_n$ 是元素 $a$ 的两个分解式, 其中 $p_i,q_j$ 均是 $R$ 中的不可约元素, 则 $m=n$, 并且存在集合 $\{1,2,\cdots,n\}$ 上的一个置换 $\sigma$, 使得 $p_i\sim q_{\sigma(i)}\ (1\leqslant i\leqslant n)$.

\textbf{(5)} 主理想整环均是 UFD. 若 $F$ 为域, 则多项式环 $F[x]$ 是主理想整环, 从而也是 UFD.

\textbf{(6)（Gauss）} 若 $R$ 是 UFD, 则多项式环 $R[x]$（从而 $R[x_1,\cdots,x_n]$）也是 UFD.

设 $a$ 和 $b$ 是环 $R$ 中两个非零元素. $d\in R$ 叫作是 $a$ 和 $b$ 的最大公因子（表示成 $(a,b)$）, 如果:

(i) $d\mid a$ 并且 $d\mid b$;

(ii) 若 $c\mid a$, $c\mid b$, $c\in R$, 则 $c\mid d$.

类似地可以定义 $a$ 和 $b$ 的最小公倍元（表示成 $[a,b]$）, 以及多个元素的最大公因子 $(a_1,a_2,\cdots,a_n)$ 和最小公倍元 $[a_1,a_2,\cdots,a_n]$. 如果 $d=(a,b)$, 则与 $d$ 相伴的每个元素均是 $a$ 和 $b$ 的最大公因子, 对于最小公倍元也有类似的结论. 如果 $R$ 是 UFD, 则任意 $n$ 个非零元素 $a_1,\cdots,a_n$ 均有最大公因子和最小公倍元. 而当 $R$ 是主理想整环时, 我们有
\[
  (a_1,\cdots,a_n)\ \text{为主理想} \Leftrightarrow
  (a_1)\cap(a_2)\cap\cdots\cap(a_n)=(d)
\]
（由元素 $a_1,\cdots,a_n$ 生成的主理想）, 其中 $d$ 是 $a_1,\cdots,a_n$ 的最大公因子或最小公倍元.

\textbf{(7)（Eisenstein 不可约判别法）} 设 $D$ 为 UFD, $F$ 是 $D$ 的商域, $f(x)=a_nx^n+\cdots+a_0\in D[x]$, $\deg f(x)\geqslant 1$, $p$ 是 $D$ 中一个不可约元素, 并且 $p\nmid a_n$, $p\mid a_i\ (0\leqslant i<n)$, $p^2\nmid a_0$, 则

(a) $f(x)$ 是 $F[x]$ 中的不可约元素;

进而,

(b) 如果在 $D$ 中还有 $(a_0,\cdots,a_n)=1$, 则 $f(x)$ 也是 $D[x]$ 中的不可约元素.

\section*{3. 域的扩张, 域的伽罗瓦理论}

设 $K$ 和 $F$ 均为域. 如果 $K$ 是 $F$ 的子域, 则称 $F$ 是 $K$ 的扩域或（域的）扩张, 并且常常把这样一对域表示成 $F/K$. 设 $F$ 是域而 $X$ 是 $F$ 的一个子集合, 则 $F$ 中包含 $X$ 的最小子域叫作是 $F$ 中由集合 $X$ 生成的子域. 如果 $K$ 是 $F$ 的子域而 $X$ 是 $F$ 的子集合, 则 $F$ 中由集合 $K\cup X$ 生成的子域也叫作是 $X$ 在域 $K$ 上生成的域, 并且表示成 $K(X)$, 而由集合 $K\cup X$ 生成的子环表示成 $K[X]$. 如果 $X=\{u_1,\cdots,u_s\}$, 则 $K(X)$ 和 $K[X]$ 也分别表示成 $K(u_1,\cdots,u_s)$ 和 $K[u_1,\cdots,u_s]$, 并且称域 $K(u_1,\cdots,u_s)$ 是 $K$ 的有限生成扩张, 特别当 $s=1$ 时, 域 $K(u)$ 叫作是域 $K$ 的单扩张. 如果 $L$ 和 $M$ 均是域 $F$ 的子域, 我们将 $F$ 中由 $L\cup M$ 生成的子域叫作是域 $L$ 和 $M$ 的合成, 并且表示成 $LM$. 于是 $LM=L(M)=M(L)=ML$. 类似地可以定义多个域的合成.

设 $F/K$ 是域的扩张, 则 $F$ 是域 $K$ 上的向量空间, 我们以 $[F:K]$ 表示向量空间 $F$ 在 $K$ 上的维数, 并且称作是扩张 $F/K$ 的次数. 当 $[F:K]$ 有限时, 称 $F/K$ 为有限（次）扩张, 否则叫作是无限（次）扩张.

\textbf{(8)} 设 $F/E$ 和 $E/K$ 均是域的扩张, 则
\[
  [F:K]=[F:E][E:K].
\]
\textbf{(9)} 设 $F$ 为 $K$ 的扩域, $X=\{x_1,x_2,\cdots\}$, 则

(i) $K(X)=\{f(x_1,\cdots,x_s)/g(x_1,\cdots,x_s)\mid f,g\in K[x_1,\cdots,x_s],\ g\neq 0,\ s\geqslant 1\}$;

(ii) $K[X]=\{f(x_1,\cdots,x_s)\mid f\in K[x_1,\cdots,x_s],\ s\geqslant 1\}$;

(iii) 对于每个元素 $\alpha\in K[X]$（或 $K(X)$）, 均存在 $X$ 的有限子集合 $X'$, 使得 $\alpha\in K[X']$（或 $K(X')$）.

设 $\alpha$ 为域 $F$ 中元素, $f(x)\in F[x]$. 如果 $f(\alpha)=0$, 则称 $\alpha$ 为多项式 $f(x)$ 的根. 根据余数定理, 这等价于在多项式环 $F[x]$ 中 $(x-\alpha)\mid f(x)$. 设 $m$ 为正整数, 使得 $(x-\alpha)^m\mid f(x)$ 而 $(x-\alpha)^{m+1}\nmid f(x)$, $m$ 叫作是根 $\alpha$ 的重数. 当 $m\geqslant 2$ 时, $\alpha$ 叫作是 $f(x)$ 的重根; 如果 $m=1$, 则 $\alpha$ 叫作是 $f(x)$ 的单根. 对于 $f(x)=\sum a_ix^i\in F[x]$, 我们将多项式 $f'(x)=\sum ia_ix^{i-1}\in F[x]$ 叫作是 $f(x)$ 的形式微商.

\textbf{(10)} 设 $F$ 为域, $f(x)\in F[x]$, $\deg f=n$:

(a) $f(x)$ 在域 $F$ 中根的个数（重根按重数计算）至多为 $n$;

(b) $f(x)$ 在 $F$ 的某个扩域中有重根 $\Leftrightarrow f(x)$ 和 $f'(x)$ 不互素;

(c) 如果 $f(x)$ 是 $F[x]$ 中的不可约多项式, 则当 $F$ 为特征零域时, $f$ 没有重根; 而当 $F$ 为特征 $p$（$p$ 为素数）域时, $f(x)$ 在 $F$ 的某个扩域中有重根的充要条件是它为 $x^{p^m}$ 的多项式, 即存在 $g(x)\in F[x]$, 使得 $f(x)=g(x^{p^m})$.

设 $F$ 是 $K$ 的扩域, 元素 $\alpha\in F$ 叫作是 $K$ 上的代数元素, 是指存在某个非零多项式 $f(x)\in K[x]$, 使得 $f(\alpha)=0$. 反之, 如果 $\alpha$ 不是 $K[x]$ 中任何非零多项式的根, 则称 $\alpha$ 是 $K$ 上的超越元素. 如果 $F$ 中每个元素均是 $K$ 上的代数元素, 则称 $F$ 是 $K$ 的代数扩张; 反之, 如果 $F$ 中至少有一个元素在 $K$ 上是超越的, 则称 $F$ 是 $K$ 的超越扩张. 域 $K$ 上的多项式环 $K[x_1,\cdots,x_n]$ 是整环, 它的商域 $K(x_1,\cdots,x_n)$ 叫作是域 $K$ 上关于未定元 $x_1,\cdots,x_n$ 的有理函数域.

\textbf{(11)} 设 $F$ 为 $K$ 的扩域, $\alpha\in F$ 是 $K$ 上的超越元素, 则

(a) 存在域的同构 $K(\alpha)\cong K(u)$, 使得 $\alpha$ 在 $K$ 上的限制 $\sigma|_K$ 是域 $K$ 的恒等自同构;

(b) $K(\alpha)/K$ 是无限（超越）扩张.

\textbf{(12)} 设 $F$ 为 $K$ 的扩域, $\alpha\in F$ 是 $K$ 上的代数元素, 则

(a) $K[x]$ 中存在惟一的不可约首 $1$（即最高项系数是 $1$）多项式 $f(x)\in K[x]$（$\deg f=n\geqslant 1$）, 使得 $f(\alpha)=0$;

(b) $K(\alpha)\cong K[x]/(f(x))$;

(c) $1,\alpha,\cdots,\alpha^{n-1}$ 是向量空间 $K(\alpha)$ 的一组 $K$-基;

(d) $K(\alpha)/K$ 是（有限）代数扩张.

\textbf{注记} (a) (12) 中所说的多项式 $f(x)$ 叫作是代数元素 $\alpha$ 在域 $K$ 上的极小多项式. 这是因为它有如下性质: (i) $f(\alpha)=0$; (ii) $g(x)\in K[x]$, $g(\alpha)=0\Rightarrow f(x)\mid g(x)$.

(b) 令 $\deg f=n=[K(\alpha):K]$, 我们也称 $\alpha$ 是 $K$ 上的 $n$ 次代数元素. (12) 表明: 这时域 $K(\alpha)$ 中每个元素均可惟一地表示成
\[
  a_0+a_1\alpha+\cdots+a_{n-1}\alpha^{n-1}\qquad(a_i\in K).
\]

\textbf{(13)} 设 $K$ 为域, $f(x)$ 是 $K[x]$ 中任一多项式, $\deg f\geqslant 1$, 则存在 $K$ 的单扩张 $F=K(u)$, 使得 $f(u)=0$. 进而, 如果 $f(x)$ 是 $K[x]$ 中的不可约多项式, 并且 $F'=K(u')$ 是 $K$ 的另一个单扩张, 使得 $f(u')=0$, 则存在域的同构 $\sigma\colon F\to F'$, 使得 $\sigma(u)=u'$ 并且 $\sigma$ 是域 $K$ 的恒等自同构.

\textbf{注记} 我们将 (13) 中的域 $K(u)$ 称作是将 $f(x)$ 的根添加到 $K$ 上而得到的域. 如果作多次的添加, 我们就可找到 $K$ 的某个扩域 $M$, 使得 $f(x)$ 的全部根都在 $M$ 中, 即 $f(x)$ 在 $M[x]$ 中分解成一次多项式的乘积. $K$ 的满足这种性质的最小扩域叫作是 $f(x)$ 在 $K$ 上的分裂域. 换句话说, 域 $M$ 叫作是多项式 $f(x)\in K[x]$ 在 $K$ 上的分裂域, 是指: $M=K(\alpha_1,\cdots,\alpha_n)$, $\alpha_i\in M$, $n=\deg f$, 并且 $f(x)=a\prod_{i=1}^{n}(x-\alpha_i)$. $f(x)$ 在 $K$ 上的任意两个分裂域是彼此同构的.

设 $K$ 是域. 如果每个在 $K$ 上代数的元素均属于 $K$, 则称 $K$ 是代数封闭域. 例如, 我们在本附录的最后要证明复数域 $\C$ 是代数封闭域. 对于任意域 $K$, 设 $\Omega$ 是将 $K$ 上代数的全部元素添加到 $K$ 上而得到的域, 则 $\Omega$ 是代数封闭域, $\Omega$ 叫作是域 $K$ 的代数闭包. $\Omega$ 的代数封闭性是基于:

\textbf{(14)} (a) 若 $\alpha$ 是 $F$ 上的代数元素而 $F/K$ 是域的代数扩张, 则 $\alpha$ 也是 $K$ 上的代数元素;

(b) 若 $F/M$, $M/K$ 均是域的代数扩张, 则 $F/K$ 也是代数扩张.

有理数域 $\Q$ 的有限（次）扩域（从而是代数扩张）$K$ 叫作是代数数域, 简称作数域, 这是代数数论的基本研究对象. 由于复数域 $\C$ 是 $\Q$ 的代数封闭扩域, 我们可以将所有的数域均看成是 $\C$ 的子域.

现在叙述数域的伽罗瓦扩张理论. 数域 $K$ 到 $\C$ 中的每个单同态也叫作是 $K$ 到 $\C$ 中的嵌入. 设 $L/K$ 是数域的扩张, $\sigma\colon L\to\C$ 为嵌入, 如果 $\sigma(K)=K$ 并且 $\sigma$ 到 $K$ 是域 $K$ 的恒等自同构, 则称 $\sigma$ 为 $L$ 到 $\C$ 中的一个 $K$-嵌入. 类似地, 设 $L_1/K$, $L_2/K$ 均是数域的扩张, 如果 $\sigma\colon L_1\to L_2$ 是域同构并且 $\sigma$ 到 $K$ 是域 $K$ 的恒等自同构, 则称 $\sigma$ 是 $K$-同构. 特别当 $L_2=L_1$ 时, 则称 $\sigma$ 为 $L_1$ 的 $K$-自同构. $L$ 的 $K$-自同构全体显然形成群, 叫作是扩张 $L/K$ 的伽罗瓦群, 表示成 $\Gal(L/K)$（注意对于 $\sigma,\tau\in\Gal(L/K)$, $\sigma\tau$ 定义为: $(\sigma\tau)(\alpha)=\sigma(\tau(\alpha))$, $\alpha\in L$）. 可以证明 $|\Gal(L/K)|\leqslant[L:K]$.

设 $L/K$ 是数域的扩张, $\sigma\colon L\to\C$ 为一个 $K$-嵌入. 对于 $\alpha\in L$, $\sigma(\alpha)\in\C$ 叫作是元素 $\alpha$ 的 $K$-共轭元素. 设 $f(x)\in K[x]$ 是元素 $\alpha$ 在 $K$ 上的极小多项式, 则 $f(x)$ 在 $\C$ 中的全部根恰好是 $\alpha$ 的全部 $K$-共轭元素. 类似地, $\sigma(L)$ 叫作是域 $L$ 的 $K$-共轭域. 数域扩张 $L/K$ 叫作是伽罗瓦扩张（或叫正规扩张）, 如果 $L$ 是 $K$-自共轭域. 也就是说, 对于每个 $K$-嵌入 $\sigma\colon L\to\C$, 均有 $\sigma(L)=L$.

\textbf{(15)} 设 $L/K$ 是数域的扩张, 则以下几条彼此等价:

(a) $L/K$ 是伽罗瓦扩张;

(b) 对于 $L$ 中每个元素 $\alpha$ 和每个 $K$-嵌入 $\sigma\colon L\to\C$, 均有 $\sigma(\alpha)\in L$;

(c) $L$ 是某个多项式 $f(x)\in K[x]$ 在 $K$ 上的分裂域.

设 $F/K$ 是数域的伽罗瓦扩张, 则对于 $F/K$ 的每个中间域 $M$（$K\leqslant M\leqslant F$）, $F/M$ 也是伽罗瓦扩张, 并且 $\Gal(F/M)$ 是 $\Gal(F/K)$ 的子群. 另一方面, 对于 $\Gal(F/K)$ 的每个子群 $H$, 令
\[
  F^H=\{\alpha\in F\mid\sigma(\alpha)=\alpha\ \text{对每个 } \sigma\in H\},
\]
这是 $F/K$ 的中间域, 叫作是 $F$ 的 $H$-固定子域.

\textbf{(16)（伽罗瓦扩张基本定理）} 设 $F/K$ 是数域的伽罗瓦扩张. 以 $\mathcal{M}$ 表示 $F/K$ 的全部中间域所组成的集合, 以 $\mathcal{H}$ 表示 $\Gal(F/K)$ 的全部子群组成的集合. 定义映射 $\varphi\colon\mathcal{M}\to\mathcal{H}$, $\varphi(M)=\Gal(F/M)$; $\psi\colon\mathcal{H}\to\mathcal{M}$, $\psi(H)=F^H$. 则

(a) 对于每个 $M\in\mathcal{M}$, $\psi(\varphi(M))=M$; 对于每个 $H\in\mathcal{H}$, $\varphi(\psi(H))=H$. 从而 $\varphi$ 和 $\psi$ 给出集合 $\mathcal{M}$ 与 $\mathcal{H}$ 之间的一一对应, 并且

(b) 设 $H_1,H_2\in\mathcal{H}$, $M_1,M_2\in\mathcal{M}$, 则 $M_1\subseteq M_2\Leftrightarrow\varphi(M_1)\supseteq\varphi(M_2)$. 进而 $\varphi(M_1\cap M_2)=\varphi(M_1)$ 和 $\varphi(M_2)$ 生成的群, $\varphi(M_1M_2)=\varphi(M_1)\cap\varphi(M_2)$;

(c) 设 $M\in\mathcal{M}$, 则 $M/K$ 为伽罗瓦扩张 $\Leftrightarrow\varphi(M)=\Gal(F/M)$ 是 $\Gal(F/K)$ 的正规子群. 并且在这个时候,
\[
  \Gal(M/K)\cong\Gal(F/K)/\Gal(F/M).
\]

\section*{4. 有限域}

\textbf{(17)} (a) 每个有限域 $F$ 均是 $q$ 元域, 其中 $q=p^n$, $n\geqslant 1$, $p$ 为素数并且是域 $F$ 的特征. $F$ 有一个 $p$ 元子域 $\F_p$, 并且 $F/\F_p$ 是 $n$ 次扩张.

(b) $q$ 元域 $F$ 的加法群是 $n$ 个 $p$ 阶循环群的直和, 而 $F^{\times}=F-\{0\}$ 是 $q-1$ 阶乘法循环群. 设 $\alpha$ 是乘法循环群 $F^{\times}$ 的一个生成元, 则 $F=\{0,1,\alpha,\alpha^2,\cdots,\alpha^{q-2}\}$, 从而 $\F_p(\alpha)=F$ 是单扩张.

(c) 固定 $\F_p$ 的一个代数闭包 $\overline{\F}_p$, 并且设 $q=p^n$, 则 $q$ 元域 $F_q$ 即是多项式 $x^{p^n}-x\in\F_p[x]$ 在 $\F_p$ 上的分裂域. 因此对于每个正整数 $n$, $\overline{\F}_p$ 中均存在惟一的 $p^n$ 元子域.

(d) 阶数相同的两个有限域彼此同构.

(e) 设 $F$ 为 $q=p^n$ 元域, 则映射 $\varphi\colon F\to F$, $\varphi(\alpha)=\alpha^p$（$\alpha\in F$）是域 $F$ 的 $n$ 阶自同构, 称作是域 $F$ 的 Frobenius 自同构. 并且 $F$ 的自同构群 $\Aut(F)$ 即是由 $\varphi$ 生成的 $n$ 阶循环群.

(f) 我们以 $\F_q$ 表示 $q$ 元域（它不计同构是惟一的）. 设 $F_{p^n}\subseteq\overline{\F}_p$, 则 $\F_{p^m}\subseteq\F_{p^n}\Leftrightarrow m\mid n$. 在这个时候, $[\F_{p^n}:\F_{p^m}]=\frac{n}{m}$, 并且 $\F_{p^n}$ 的 $\F_{p^m}$-自同构群
\[
  \Gal(\F_{p^n}/\F_{p^m})=\{\sigma\in\Aut(\F_{p^n})\mid\sigma(a)=a\ \text{对每个 } a\in\F_{p^m}\}
\]
是由 $\varphi^m$ 生成的 $\frac{n}{m}$ 阶循环群, 从而 $\F_{p^n}/\F_{p^m}$ 是伽罗瓦扩张. 我们将 $\varphi_{n,m}\colon F_{p^n}\to F_{p^n}$, $\varphi_{n,m}(\alpha)=\alpha^{p^m}$（$\alpha\in F_{p^n}$）称作是扩张 $F_{p^n}/F_{p^m}$ 的 Frobenius 自同构. 另一方面, 对于 $\Gal(\F_{p^n}/\F_p)$ 中惟一的 $\frac{n}{m}$ 阶循环子群 $\langle\varphi^m\rangle$, 它的固定子域恰好是 $\F_{p^m}$, 即 $F_{p^n}^{\langle\varphi^m\rangle}=\F_{p^m}$.

\section*{5. 代数基本定理}

\textbf{(18)} 复数域 $\C$ 是代数封闭域. 这相当于说: 每个复系数多项式 $f(x)$（$\deg f(x)\geqslant 1$）必有复根, 并且熟知它也等价于 $n$ 次复系数多项式恰好有 $n$ 个复根（计算重数）. 有人统计, 代数基本定理共有近两百个证明. 更有趣的是, 所有的证明都要利用数学分析中的某些事实, 我们这里所介绍的一个证明也不例外, 因为我们要利用

\textbf{引理 1} 每个奇次实系数多项式必有实根. 大家知道, 这一事实从连续函数的中值定理得出.

此外我们还需要（与数学分析无关的）

\textbf{引理 2} 复系数 $2$ 次多项式的根均是复根.

\textbf{证明} 设 $f(x)$ 是复系数 $2$ 次多项式, 不妨设它是首 $1$ 的, 即 $f(x)=x^2+\alpha x+\beta\in\C[x]$. 由于 $f(x)=g(x/\gamma)$ 型变换, 我们只需证明每个多项式 $g(z)=z^2-\gamma$ 的根均是复根即可. 令 $\gamma=a+bi$, $a,b\in\R$. 取
\[
  \alpha_0=\begin{cases}
    \sqrt{\frac{a+\sqrt{a^2+b^2}}{2}}, & \text{若 } b\geqslant 0,\\
    \sqrt{\frac{a-\sqrt{a^2+b^2}}{2}}, & \text{若 } b<0,
  \end{cases}
\]
\[
  \beta_0=\begin{cases}
    \sqrt{\frac{-a+\sqrt{a^2+b^2}}{2}}, & \text{若 } b\geqslant 0,\\
    -\sqrt{\frac{-a+\sqrt{a^2+b^2}}{2}}, & \text{若 } b<0,
  \end{cases}
\]
直接验证即知 $\alpha_0+i\beta_0$ 和 $-\alpha_0-i\beta_0$ 是 $g(z)=z^2-\gamma$ 的两个复根.

\textbf{引理 3} 每个次数 $\geqslant 1$ 的实系数多项式 $g(x)$ 必有复根.

\textbf{证明} 设 $g(x)\in\R[x]$, $\deg g(x)=d\geqslant 1$. 令 $d=2^nq$, $2\nmid q$, $n\geqslant 0$. 我们对 $n$ 作数学归纳法. 如果 $n=0$, 则 $g(x)$ 为奇次实系数多项式, 由引理 1 可知它有实根. 下设 $n=n_0\geqslant 1$, 并且对于比 $n_0$ 小的情形引理 3 成立. 令 $\alpha_1,\cdots,\alpha_d$ 为 $g(x)$ 在 $\C$ 的适当的扩域中的全部根, 我们的目的是证明必有某个 $\alpha_i\in\C$. 为证此, 取任意一个实数 $c$, 令
\[
  \beta_{ij}=\alpha_i\alpha_j+c(\alpha_i+\alpha_j)\qquad(1\leqslant i<j\leqslant d),
\]
共有 $\binom{d+1}{2}=2^{n-1}q(d+1)$ 个, 又令
\[
  G(x)=\prod_{i<j}(x-\beta_{ij})=x^{2^{n-1}q(d+1)}+g_1(x)x^{2^{n-1}q(d+1)-1}+\cdots,
\]
不难看出, 每个 $g_k\ (1\leqslant k\leqslant 2^{n-1}q(d+1))$ 均是 $\alpha_1,\cdots,\alpha_d$ 与 $c$ 的实系数对称多项式. 由于 $\alpha_i\ (1\leqslant i\leqslant d)$ 是 $g(x)$ 的根, 而 $g(x)$ 的系数均为实数, 从而 $g_k$ 均是 $c$ 和 $g(x)$ 的系数的多项式, 即 $g_k\in\R[c]$. 从而 $G(x)\in\R[x]$. 由于 $n_0\geqslant 1$, 从而 $\deg G(x)=2^{n-1}q(d+1)$, 而 $2^{n-1}$ 中 $2$ 的幂次小于 $2^n$, 由归纳假设便知 $G(x)$ 有复根 $\beta_0$. 换句话说, 我们有 $i,j$ 和 $i',j'$（均与 $c$ 有关）, 使得
\[
  (\alpha_i+\alpha_j)+c\alpha_i\alpha_j=\beta_0
  \quad\text{并且}\quad
  (\alpha_{i'}+\alpha_{j'})+c\alpha_{i'}\alpha_{j'}=\beta_0.
\]
由于指标集合 $\{(i,j)\mid i<j\}$ 是有限的, 而实数 $c$ 可任意选取, 从而必然存在 $c\neq c'$, $c,c'\in\R$, 使得 $\beta_0(c)=\beta_0(c')$ 但对应的指标对不同. 令它们分别为 $\{\alpha_1,\alpha_2\}$ 和 $\{\alpha_3,\alpha_4\}$, 于是得到
\[
  \alpha_1+\alpha_2+c\alpha_1\alpha_2=\alpha_3+\alpha_4+c\alpha_3\alpha_4=\beta_0,
\]
\[
  \alpha_1+\alpha_2+c'\alpha_1\alpha_2=\alpha_3+\alpha_4+c'\alpha_3\alpha_4=\beta_0.
\]
由此及 $c\neq c'$ 可知 $\alpha_1\alpha_2=\alpha_3\alpha_4$ 且 $\alpha_1+\alpha_2=\alpha_3+\alpha_4$, 于是 $h(x)=(x-\alpha_1)(x-\alpha_2)=(x-\alpha_3)(x-\alpha_4)$ 是实系数 $2$ 次多项式. 根据引理 2, $h(x)$ 的两个根 $\alpha_1$ 和 $\alpha_2$ 均是复根, 这就证明了存在某个 $1\leqslant i\leqslant d$, 使得 $\alpha_i\in\C$.

最后我们来证代数基本定理: 设 $f(x)\in\C[x]$, $\deg f\geqslant 1$. 记 $f(x)=\sum a_ix^i$, 其中 $\bar{a}_i$ 表示 $a_i$ 的共轭复数. 令 $g(x)=f(x)\overline{f(x)}$, 易知 $g(x)\in\R[x]$. 根据引理 3, 存在 $\alpha\in\C$, 使得 $g(\alpha)=0$, 于是 $\alpha$ 为 $f(x)$ 或者 $\overline{f(x)}$ 的根. 如果 $f(\alpha)=0$ 则证毕; 如果 $\overline{f(\alpha)}=0$, 则复数 $\bar{\alpha}$ 就是 $f(x)$ 的根. 这就证明了代数基本定理.
